\documentclass[11pt,a4paper]{article}

% ------------------ PAQUETES ------------------
\usepackage[utf8]{inputenc}
\usepackage[T1]{fontenc}
\usepackage[spanish]{babel}
\usepackage{geometry}
\usepackage{pgfgantt}
\usepackage{float}
\usepackage{array}
\usepackage{colortbl} % Para colorear celdas
\usepackage{xcolor}

\geometry{margin=2cm} % Márgenes más estrechos para aprovechar espacio

% ------------------ ESTILOS ------------------
\definecolor{headergray}{gray}{0.9}
\renewcommand{\arraystretch}{1.4}

\begin{document}

% ------------------ CABECERA SIMPLIFICADA ------------------
\begin{center}
    {\Large \textbf{Planificación TFG: Resumen concentrado}}\\[5pt]
    {\large Agente LLM autónomo para auditorías de ciberseguridad}
    
    \vspace{0.5cm}
    \textbf{Alumno:} Alejandro Barrena Millán \quad | \quad \textbf{Inicio:} 17 Nov 2025 \quad | \quad \textbf{Fin:} Feb 2026
\end{center}

\hrule
\vspace{0.5cm}

% ============================================================
\section*{1. Objetivo y Metodología}

\begin{itemize}
    \item \textbf{Objetivo:} Desarrollar un prototipo funcional de agente basado en LLM capaz de ejecutar comandos de auditoría en un entorno aislado.
    \item \textbf{Metodología:} Scrum adaptado (Ciclos incrementales de 1 a 3 semanas).
\end{itemize}

% ============================================================
\section*{2. Hoja de Ruta y Entregables (Roadmap)}

Resumen de hitos clave y fechas de calendario real.

\begin{table}[H]
\centering
\begin{tabular}{|c|l|p{8.5cm}|}
\hline
\rowcolor{headergray} 
\textbf{Sprint} & \textbf{Fechas (Inicio)} & \textbf{Entregable Principal / Objetivo} \\ \hline

\textbf{S1} & 17 Nov & \textbf{Definición:} Requisitos, GitHub configurado y estado del arte. \\ \hline

\textbf{S2} & 24 Nov - 07 Dic & \textbf{Infraestructura:} Sandbox (Docker) operativo y API de control lista. \\ \hline

\textbf{S3} & 08 Dic - 28 Dic & \textbf{Integración IA:} El LLM se conecta al entorno y ejecuta comandos simples. \\ \hline

\textbf{S4} & 29 Dic - 11 Ene & \textbf{MVP:} Agente autónomo realizando auditorías básicas (puertos/servicios). \\ \hline

\textbf{S5} & 12 Ene - 18 Ene & \textbf{Validación:} Pruebas end-to-end, corrección de errores y métricas. \\ \hline

\textbf{S6} & 19 Ene - 01 Feb & \textbf{Memoria:} Redacción completa del documento académico. \\ \hline

\textbf{S7} & 02 Feb & \textbf{Defensa:} Entrega final y preparación de presentación. \\ \hline
\end{tabular}
\end{table}

% ============================================================
\section*{3. Cronograma Visual}

Vista rápida de la carga de trabajo por semanas.

\begin{figure}[H]
\centering
\begin{ganttchart}[
    hgrid,
    vgrid,
    x unit=1.2cm,
    y unit chart=0.5cm,
    bar/.append style={fill=blue!40, draw=none},
    bar height=0.6,
    title label font=\bfseries\footnotesize,
    bar label font=\footnotesize
]{1}{12}

\gantttitle{Semanas (Nov 2025 - Feb 2026)}{12} \\
\gantttitlelist{1,...,12}{1} \\

\ganttbar{S1: Requisitos}{1}{1} \\
\ganttbar{S2: Entorno}{2}{3} \\
\ganttbar{S3: Integración LLM}{4}{6} \\ 
\ganttbar{S4: MVP Agente}{7}{8} \\
\ganttbar{S5: Pruebas}{9}{9} \\
\ganttbar{S6: Memoria}{10}{11} \\
\ganttbar{S7: Defensa}{12}{12}      

\end{ganttchart}
\end{figure}

\end{document}